% =============================================================================
%  IIIT RAICHUR -- B.Tech PROJECT THESIS TEMPLATE
%  Department of Computer Science and Engineering
%  -----------------------------------------------------------------------------
%  HOW TO USE
%   1. Edit the placeholder commands in the "USER INPUTS" block below.
%   2. Place your college logo at: images/iiitr_logo.png
%   3. Replace the lorem ipsum / dummy text with your own project content.
%   4. Add your figures to the images/ folder.
%   5. Compile:
%        pdflatex thesis.tex
%        pdflatex thesis.tex
%      (or simply: latexmk -pdf thesis.tex)
%
%  This template is domain-agnostic. The dummy text is generic so it works for
%  any kind of project -- software, hardware, embedded, ML, networking, etc.
% =============================================================================

\documentclass[12pt, a4paper, oneside]{report}

% ---------------------- PACKAGES ---------------------------------------------
\usepackage[utf8]{inputenc}
\usepackage[T1]{fontenc}
\usepackage{lmodern}
\usepackage[margin=1in, top=1.2in, bottom=1.2in]{geometry}
\usepackage{graphicx}
\usepackage{setspace}
\usepackage{titlesec}
\usepackage{tocloft}
\usepackage{fancyhdr}
\usepackage{booktabs}
\usepackage{longtable}
\usepackage{array}
\usepackage{tabularx}
\usepackage{multirow}
\usepackage{caption}
\usepackage{subcaption}
\usepackage{amsmath, amssymb}
\usepackage{xcolor}
\usepackage[hidelinks, breaklinks=true]{hyperref}
\usepackage{url}
\usepackage{enumitem}
\usepackage{float}
\usepackage{lipsum}    % provides \lipsum for lorem ipsum dummy text

% ---------------------- USER INPUTS (EDIT THESE) -----------------------------
% Plain title (single line) -- used in declaration, certificate, abstract.
\newcommand{\thesisTitle}{Title of Your Project as a Single Line}
% Display title (with line breaks) -- used only on the title page.
% Add \\ wherever you want a line break for visual layout.
\newcommand{\thesisTitleDisplay}{TITLE OF YOUR PROJECT \\ IN TWO OR THREE LINES \\ AS NEEDED}
\newcommand{\studentName}{Your Full Name}
\newcommand{\rollNumber}{CS21B10XX}
\newcommand{\supervisorName}{Dr.~Supervisor Name}
\newcommand{\departmentName}{Computer Science and Engineering}
\newcommand{\submissionMonthYear}{May 2025}
\newcommand{\submissionLocation}{Raichur -- 584135}
\newcommand{\degreeName}{Bachelor of Technology}
\newcommand{\logoPath}{images/iiitr_logo.png}

% ---------------------- FORMATTING -------------------------------------------
\onehalfspacing

% Chapter title formatting (matches "Chapter N" on top, title below)
\titleformat{\chapter}[display]
  {\normalfont\huge\bfseries}{\chaptertitlename\ \thechapter}{20pt}{\Huge}
\titlespacing*{\chapter}{0pt}{-30pt}{30pt}

% Page style -- page number at bottom centre, no header rule
\pagestyle{fancy}
\fancyhf{}
\fancyfoot[C]{\thepage}
\renewcommand{\headrulewidth}{0pt}
\renewcommand{\footrulewidth}{0pt}

\fancypagestyle{plain}{%
  \fancyhf{}%
  \fancyfoot[C]{\thepage}%
  \renewcommand{\headrulewidth}{0pt}%
  \renewcommand{\footrulewidth}{0pt}%
}

% TOC formatting
\renewcommand{\cftchapfont}{\bfseries}
\renewcommand{\cftchappagefont}{\bfseries}

% Helper command for centered front-matter headings (Declaration, Certificate, Abstract)
\newcommand{\frontmatterheading}[1]{%
  \begin{center}\Large\bfseries #1\end{center}
  \vspace{0.5cm}
}

% Paragraph spacing
\setlength{\parskip}{0.5em}
\setlength{\parindent}{1.5em}

% =============================================================================
\begin{document}

% ---------------------- TITLE PAGE -------------------------------------------
\begin{titlepage}
    \centering
    \thispagestyle{empty}
    \vspace*{0.2cm}

    {\Large\bfseries \thesisTitleDisplay \par}
    \vspace{0.8cm}

    {\normalsize A Project Report Submitted \par}
    {\normalsize in Partial Fulfilment of the Requirements \par}
    {\normalsize for the Degree of \par}
    \vspace{0.5cm}

    {\large\bfseries \MakeUppercase{\degreeName} \par}
    \vspace{0.3cm}

    {\normalsize in \par}
    \vspace{0.2cm}

    {\large\bfseries \MakeUppercase{\departmentName} \par}
    \vspace{0.5cm}

    {\normalsize\itshape by \par}
    \vspace{0.3cm}

    {\normalsize\bfseries \studentName{} (Roll No.~\rollNumber) \par}
    \vspace{0.5cm}

    % College logo
    \includegraphics[width=1\textwidth]{\logoPath}\par
    \vspace{0.5cm}

    {\normalsize\itshape to \par}
    \vspace{0.2cm}

    {\normalsize\bfseries DEPARTMENT OF \MakeUppercase{\departmentName} \par}
    {\normalsize\bfseries INDIAN INSTITUTE OF INFORMATION TECHNOLOGY \par}
    {\normalsize\bfseries RAICHUR-584135, INDIA \par}
    \vspace{0.4cm}

    {\normalsize\itshape \submissionMonthYear \par}

\end{titlepage}

% ---------------------- FRONT MATTER (roman page numbers) --------------------
\pagenumbering{roman}
\setcounter{page}{2}

% ---------------------- DECLARATION ------------------------------------------
\clearpage
\thispagestyle{plain}
\frontmatterheading{DECLARATION}
\addcontentsline{toc}{chapter}{Declaration}

\noindent I \textbf{\MakeUppercase{\studentName}} (Roll Number: \rollNumber) hereby
declare that, this report entitled ``\textit{\thesisTitle}'' submitted to
Indian Institute of Information Technology Raichur towards partial requirement
of \degreeName{} in \departmentName{} is an original work carried out by me
under the supervision of Mentor, \supervisorName{} and has not formed the
basis for the award of any degree or diploma, in this or any other institution
or university. I have sincerely tried to uphold the academic ethics and
honesty. Whenever an external information or statement or result is used then,
that have been duly acknowledged and cited.

\vspace{2.5cm}
\noindent
\begin{minipage}[t]{0.5\textwidth}
\submissionLocation \\
\submissionMonthYear
\end{minipage}
\hfill
\begin{minipage}[t]{0.4\textwidth}
\raggedleft
\studentName \\
Roll No.~\rollNumber
\end{minipage}

% ---------------------- CERTIFICATE ------------------------------------------
\clearpage
\thispagestyle{plain}
\frontmatterheading{CERTIFICATE}
\addcontentsline{toc}{chapter}{Certificate}

\noindent This is to certify that the work contained in this project report
entitled ``\textit{\thesisTitle}'' submitted by \textbf{\studentName} (Roll
No: \rollNumber) to the Indian Institute of Information Technology Raichur
towards partial requirement of \degreeName{} in \departmentName{} has been
carried out by him/her under my supervision and that it has not been submitted
elsewhere for the award of any degree.

\vspace{3cm}
\noindent
\begin{minipage}[t]{0.5\textwidth}
\submissionLocation \\
\submissionMonthYear
\end{minipage}
\hfill
\begin{minipage}[t]{0.4\textwidth}
\raggedleft
(\supervisorName) \\
Project Supervisor
\end{minipage}

% ---------------------- ABSTRACT ---------------------------------------------
\clearpage
\thispagestyle{plain}
\frontmatterheading{ABSTRACT}
\addcontentsline{toc}{chapter}{Abstract}

% Replace this dummy text with your own abstract (~250-300 words).
% A good abstract: motivates the problem, states what you built, summarises
% the methodology, and highlights the key results.
\lipsum[1]

\lipsum[2][1-5]

% ---------------------- TABLE OF CONTENTS ------------------------------------
\clearpage
\tableofcontents
\listoffigures
\addcontentsline{toc}{chapter}{List of Figures}
\listoftables
\addcontentsline{toc}{chapter}{List of Tables}

\clearpage
\pagenumbering{arabic}
\setcounter{page}{1}

% =============================================================================
% CHAPTER 1: INTRODUCTION
% =============================================================================
\chapter{Introduction}
\label{chap:introduction}

\lipsum[3][1-7]
% Replace the lorem ipsum above with your introductory paragraph -- set the
% broader context of your project, discuss the relevance of the domain, and
% motivate the structure of this chapter.

\section{Problem Statement}
\label{sec:problem-statement}

\lipsum[4][1-2]

\begin{itemize}
    \item \textbf{Issue One:} \lipsum[5][1-2]
    \item \textbf{Issue Two:} \lipsum[6][1-2]
    \item \textbf{Issue Three:} \lipsum[7][1-2]
    \item \textbf{Issue Four:} \lipsum[8][1-2]
\end{itemize}

\section{Project Significance}
\label{sec:project-significance}

\lipsum[9][1-3]

\begin{itemize}
    \item \textbf{Benefit One:} \lipsum[10][1-2]
    \item \textbf{Benefit Two:} \lipsum[11][1-2]
    \item \textbf{Benefit Three:} \lipsum[12][1-2]
    \item \textbf{Benefit Four:} \lipsum[13][1-2]
\end{itemize}

\section{Scope of Work}
\label{sec:scope}

\lipsum[14][1-3]

\begin{itemize}
    \item \textbf{Component A:} \lipsum[15][1-2]
    \item \textbf{Component B:} \lipsum[16][1-2]
    \item \textbf{Component C:} \lipsum[17][1-2]
\end{itemize}

\lipsum[18][1-3]

\section{Objectives}
\label{sec:objectives}

The main objectives of this project are:

\begin{enumerate}
    \item \lipsum[19][1-1]
    \item \lipsum[20][1-1]
    \item \lipsum[21][1-1]
    \item \lipsum[22][1-1]
    \item \lipsum[23][1-1]
    \item \lipsum[24][1-1]
\end{enumerate}

\section{Novelty}
\label{sec:novelty}

\lipsum[25][1-3]

\begin{itemize}
    \item \textbf{Novel Aspect One:} \lipsum[26][1-2]
    \item \textbf{Novel Aspect Two:} \lipsum[27][1-2]
    \item \textbf{Novel Aspect Three:} \lipsum[28][1-2]
\end{itemize}

\lipsum[29][1-3]

% =============================================================================
% CHAPTER 2: LITERATURE SURVEY
% =============================================================================
\chapter{Literature Survey}
\label{chap:literature}

\lipsum[30][1-7]
% Replace with an overview of the literature in your domain. Identify the
% major directions / categories of prior work and the gaps your project fills.

\section{Existing Approaches -- Category A}
\label{sec:category-a}

\subsection{Approach A.1}
\lipsum[31][1-3] \cite{ref1}

\subsection{Approach A.2}
\lipsum[32][1-3] \cite{ref2}

\subsection{Approach A.3}
\lipsum[33][1-3] \cite{ref3}

\subsection{Differentiation of Our Approach}
\lipsum[34][1-2]

\begin{itemize}
    \item \lipsum[35][1-1]
    \item \lipsum[36][1-1]
\end{itemize}

\lipsum[37][1-2]

\begin{table}[H]
\centering
\caption{Summary of Category A Approaches and Our Approach}
\label{tab:cat-a-comparison}
\small
\begin{tabularx}{\textwidth}{|l|l|l|X|X|}
\hline
\textbf{Author(s)} & \textbf{Title} & \textbf{Year} & \textbf{Methodology} & \textbf{Key Findings} \\
\hline
Author A \cite{ref1} & Approach A.1 & 20XX & Methodology of approach A.1 & Strengths and limitations of approach A.1 \\
\hline
Author B \cite{ref2} & Approach A.2 & 20XX & Methodology of approach A.2 & Strengths and limitations of approach A.2 \\
\hline
Author C \cite{ref3} & Approach A.3 & 20XX & Methodology of approach A.3 & Strengths and limitations of approach A.3 \\
\hline
\textbf{Proposed} & Our project name & 20XX & Brief summary of our methodology & Brief summary of our contribution \\
\hline
\end{tabularx}
\end{table}

\section{Existing Approaches -- Category B}
\label{sec:category-b}

\lipsum[38][1-4]

\subsection{Approach B.1}
\lipsum[39][1-3] \cite{ref4}

\subsection{Approach B.2}
\lipsum[40][1-3] \cite{ref5}

\subsection{Approach B.3}
\lipsum[41][1-3] \cite{ref6}

\subsection{Differentiation of Our Approach}
\lipsum[42][1-2]

\begin{itemize}
    \item \textbf{Feature One:} \lipsum[43][1-1]
    \item \textbf{Feature Two:} \lipsum[44][1-1]
    \item \textbf{Feature Three:} \lipsum[45][1-1]
\end{itemize}

\begin{table}[H]
\centering
\caption{Comparison of Category B Approaches}
\label{tab:cat-b-comparison}
\small
\begin{tabularx}{\textwidth}{|l|X|X|}
\hline
\textbf{Tool / Method} & \textbf{Key Features} & \textbf{Limitations} \\
\hline
Approach B.1 \cite{ref4} & Key features summary & Limitations summary \\
\hline
Approach B.2 \cite{ref5} & Key features summary & Limitations summary \\
\hline
Approach B.3 \cite{ref6} & Key features summary & Limitations summary \\
\hline
\end{tabularx}
\end{table}

\section{Existing Approaches -- Category C}
\label{sec:category-c}

\lipsum[46][1-4]

\subsection{Approach C.1}
\lipsum[47][1-2] \cite{ref7}

\subsection{Approach C.2}
\lipsum[48][1-2] \cite{ref8}

\subsection{Differentiation of Our Approach}
\lipsum[49][1-3]

\begin{itemize}
    \item \textbf{Aspect One:} \lipsum[50][1-1]
    \item \textbf{Aspect Two:} \lipsum[51][1-1]
\end{itemize}

\begin{table}[H]
\centering
\caption{Comparison of Category C Approaches}
\label{tab:cat-c-comparison}
\small
\begin{tabularx}{\textwidth}{|l|X|X|X|}
\hline
\textbf{Tool / Method} & \textbf{Key Features} & \textbf{Methodology} & \textbf{Limitations} \\
\hline
Approach C.1 \cite{ref7} & Features & Methodology & Limitations \\
\hline
Approach C.2 \cite{ref8} & Features & Methodology & Limitations \\
\hline
\end{tabularx}
\end{table}

\lipsum[52][1-2]

% =============================================================================
% CHAPTER 3: METHODOLOGY
% =============================================================================
\chapter{Methodology}
\label{chap:methodology}

\lipsum[53][1-4]

\section{Problem Restatement}
\label{sec:problem-restatement}

\lipsum[54][1-3]

\begin{itemize}
    \item \lipsum[55][1-1]
    \item \lipsum[56][1-1]
    \item \lipsum[57][1-1]
\end{itemize}

\section{Overall Approach}
\label{sec:overall-approach}

\lipsum[58][1-2]

\begin{itemize}
    \item \textbf{Module One:} \lipsum[59][1-2]
    \item \textbf{Module Two:} \lipsum[60][1-2]
    \item \textbf{Module Three:} \lipsum[61][1-2]
\end{itemize}

\section{System Block Diagram / Data Model}
\label{sec:block-diagram}

The block diagram (Figure~\ref{fig:block-diagram}) shows the high-level
structure of the system and the relationships among its core entities.

\begin{figure}[H]
    \centering
    \includegraphics[width=0.9\textwidth]{images/block_diagram.png}
    \caption{System Block Diagram}
    \label{fig:block-diagram}
\end{figure}

\lipsum[62][1-3]

\section{Architecture}
\label{sec:architecture}

The overall architecture is shown in Figure~\ref{fig:architecture}.
\lipsum[63][1-4]

\begin{figure}[H]
    \centering
    \includegraphics[width=0.9\textwidth]{images/architecture.png}
    \caption{System Architecture}
    \label{fig:architecture}
\end{figure}

\section{Module One}
\label{sec:module-one}

\lipsum[64][1-3]

\begin{figure}[H]
    \centering
    \includegraphics[width=0.9\textwidth]{images/module1_flow.png}
    \caption{Module One -- Process Flow}
    \label{fig:module1}
\end{figure}

\section{Module Two}
\label{sec:module-two}

\lipsum[65][1-3]

\begin{figure}[H]
    \centering
    \includegraphics[width=0.9\textwidth]{images/module2_flow.png}
    \caption{Module Two -- Process Flow}
    \label{fig:module2}
\end{figure}

\section{Module Three}
\label{sec:module-three}

\lipsum[66][1-3]

\begin{figure}[H]
    \centering
    \includegraphics[width=0.9\textwidth]{images/module3_flow.png}
    \caption{Module Three -- Process Flow}
    \label{fig:module3}
\end{figure}

\section{Module Four}
\label{sec:module-four}

\lipsum[67][1-3]

\begin{figure}[H]
    \centering
    \includegraphics[width=0.9\textwidth]{images/module4_flow.png}
    \caption{Module Four -- Process Flow}
    \label{fig:module4}
\end{figure}

% =============================================================================
% CHAPTER 4: CHALLENGES AND SOLUTIONS
% =============================================================================
\chapter{Challenges and Solutions}
\label{chap:challenges}

\lipsum[68][1-5]

A summary of the major issues faced and their corresponding solutions is
provided in Table~\ref{tab:challenges}.

\begin{table}[H]
\centering
\caption{Major Implementation Challenges and Their Solutions}
\label{tab:challenges}
\begin{tabularx}{\textwidth}{|X|X|}
\hline
\textbf{Challenge} & \textbf{Solution} \\
\hline
Description of the first challenge encountered during implementation
& Description of how the first challenge was resolved \\
\hline
Description of the second challenge encountered during implementation
& Description of how the second challenge was resolved \\
\hline
Description of the third challenge encountered during implementation
& Description of how the third challenge was resolved \\
\hline
\end{tabularx}
\end{table}

% =============================================================================
% CHAPTER 5: RESULTS AND DISCUSSION
% =============================================================================
\chapter{Results and Discussion}
\label{chap:results}

\lipsum[69][1-3]

\section{Result Set One}
\label{sec:result-set-1}

\lipsum[70][1-2]

\subsection{Output 1.1}
\lipsum[71][1-2]
(See Figure~\ref{fig:result1}.)

\begin{figure}[H]
    \centering
    \includegraphics[width=0.9\textwidth]{images/result1.png}
    \caption{Caption describing output 1.1}
    \label{fig:result1}
\end{figure}

\subsection{Output 1.2}
\lipsum[72][1-2]
(See Figure~\ref{fig:result2}.)

\begin{figure}[H]
    \centering
    \includegraphics[width=0.9\textwidth]{images/result2.png}
    \caption{Caption describing output 1.2}
    \label{fig:result2}
\end{figure}

\subsection{Output 1.3}
\lipsum[73][1-2]

\begin{figure}[H]
    \centering
    \includegraphics[width=0.9\textwidth]{images/result3.png}
    \caption{Caption describing output 1.3}
    \label{fig:result3}
\end{figure}

\section{Result Set Two}
\label{sec:result-set-2}

\subsection{Output 2.1}
\lipsum[74][1-2]

\begin{figure}[H]
    \centering
    \includegraphics[width=0.9\textwidth]{images/result4.png}
    \caption{Caption describing output 2.1}
    \label{fig:result4}
\end{figure}

\subsection{Output 2.2}
\lipsum[75][1-2]

\begin{figure}[H]
    \centering
    \includegraphics[width=0.9\textwidth]{images/result5.png}
    \caption{Caption describing output 2.2}
    \label{fig:result5}
\end{figure}

% Add more result subsections as needed.

% =============================================================================
% CHAPTER 6: TOOLS, TECHNOLOGIES, AND PARAMETERS
% =============================================================================
\chapter{Tools, Technologies, and Parameters}
\label{chap:tools}

\lipsum[76][1-5]

\section{Tools and Technologies}
\label{sec:tools}

\begin{itemize}
    \item \textbf{Programming Languages / Frameworks:} List the languages and
          frameworks used (e.g.~C/C++, Python, MATLAB, Verilog, Java, etc.).
    \item \textbf{Hardware / Platform:} List any hardware components, boards,
          sensors, or platforms (or write ``Not applicable'' for pure-software
          projects).
    \item \textbf{Libraries / Toolkits:} List the major libraries or toolkits
          relied upon (e.g.~scientific computing, simulation, ML, signal
          processing libraries).
    \item \textbf{Datasets / External Resources:} List datasets, APIs, or
          external resources used.
    \item \textbf{Development \& Build Tools:} List IDEs, build systems,
          version control, simulation tools, etc.
\end{itemize}

\section{Key Parameters and Settings}
\label{sec:parameters}

\begin{table}[H]
\centering
\caption{Important Parameters Used in Implementation}
\label{tab:parameters}
\begin{tabularx}{\textwidth}{|l|X|}
\hline
\textbf{Parameter} & \textbf{Value / Description} \\
\hline
Parameter 1 & Value and short rationale \\
\hline
Parameter 2 & Value and short rationale \\
\hline
Parameter 3 & Value and short rationale \\
\hline
Parameter 4 & Value and short rationale \\
\hline
\end{tabularx}
\end{table}

% =============================================================================
% CHAPTER 7: CONCLUSION AND FUTURE SCOPE
% =============================================================================
\chapter{Conclusion and Future Scope}
\label{chap:conclusion}

\lipsum[77][1-5]

\noindent\textbf{Future Enhancements:}

\begin{itemize}
    \item \lipsum[78][1-1]
    \item \lipsum[79][1-1]
    \item \lipsum[80][1-1]
    \item \lipsum[81][1-1]
    \item \lipsum[82][1-1]
\end{itemize}

\lipsum[83][1-2]

% =============================================================================
% BIBLIOGRAPHY
% =============================================================================
\addcontentsline{toc}{chapter}{Bibliography}

% Option 1: Manual entries (default -- start writing references right away)
\begin{thebibliography}{99}

\bibitem{ref1}
Author Surname, Initials. (Year). \textit{Title of the work}.
Journal / Publisher / URL.

\bibitem{ref2}
Author Surname, Initials. (Year). \textit{Title of the work}.
Journal / Publisher / URL.

\bibitem{ref3}
Author Surname, Initials. (Year). \textit{Title of the work}.
Journal / Publisher / URL.

\bibitem{ref4}
Author Surname, Initials. (Year). \textit{Title of the work}.
Journal / Publisher / URL.

\bibitem{ref5}
Author Surname, Initials. (Year). \textit{Title of the work}.
Journal / Publisher / URL.

\bibitem{ref6}
Author Surname, Initials. (Year). \textit{Title of the work}.
Journal / Publisher / URL.

\bibitem{ref7}
Author Surname, Initials. (Year). \textit{Title of the work}.
Journal / Publisher / URL.

\bibitem{ref8}
Author Surname, Initials. (Year). \textit{Title of the work}.
Journal / Publisher / URL.

% Add more \bibitem entries as needed.

\end{thebibliography}

% Option 2: Use BibTeX instead. To switch:
%   1. Comment out the manual \begin{thebibliography} ... \end{thebibliography}
%      block above.
%   2. Create a references.bib file in this folder.
%   3. Uncomment the two lines below.
%   4. Compile with: pdflatex -> bibtex -> pdflatex -> pdflatex
%
% \bibliographystyle{plain}
% \bibliography{references}

\end{document}
